% Compile with: pdflatex pitch_document.tex
% DoD SBIR Phase I Technical Volume — Navy Topic NV029
% Low-Cost Moldable Metastructure Materials for Maritime Protective Applications
\documentclass[11pt,letterpaper]{article}

\usepackage[T1]{fontenc}
\usepackage{lmodern}
\usepackage[margin=1in]{geometry}
\usepackage{hyperref}
\usepackage{enumitem}
\usepackage{parskip}
\usepackage{xcolor}
\usepackage{fancyhdr}
\usepackage{graphicx}
\usepackage{booktabs}

\hypersetup{
  colorlinks=true,
  linkcolor=blue!70!black,
  urlcolor=blue!60!black,
}

\definecolor{navyblue}{HTML}{003366}

\pagestyle{fancy}
\fancyhf{}
\rhead{\textcolor{gray}{\small DoD SBIR Phase I --- Topic NV029}}
\lhead{\textcolor{gray}{\small [Company Name]}}
\rfoot{\textcolor{gray}{\small Page \thepage\ of 5}}

\title{%
  \vspace{-1.5cm}
  \textcolor{navyblue}{\Large\textbf{DoD SBIR Phase I Technical Volume}}\\[0.3cm]
  \textcolor{navyblue}{\large Topic NV029: Low-Cost Moldable Metastructure Materials\\
  for Maritime Protective Applications}\\[0.4cm]
  \textbf{AI-Driven Platform for Discovery of\\
  Multifunctional Maritime Protective Coatings}\\[0.3cm]
  \large [Company Name]\\[0.2cm]
  \normalsize Principal Investigator: [PI Name]\\
  \normalsize [Address, City, State, ZIP]\\
  \normalsize CAGE Code: [XXXXX] \quad DUNS: [XXXXXXXXX]
}
\author{}
\date{March 2026}

\begin{document}
\maketitle
\thispagestyle{fancy}

\noindent\textcolor{navyblue}{\rule{\textwidth}{1.5pt}}

%% =======================================================================
\section{Technical Approach}
\label{sec:technical_approach}
%% =======================================================================

\subsection{Problem Statement and Navy Need}

Maritime assets---ship hulls, topside structures, sensor housings, unmanned vehicle
airframes, and shore-based equipment---face a uniquely aggressive combination of
environmental stressors: broadband UV irradiance (290--400\,nm at 40--60\,W/m$^2$
in tropical theaters), continuous salt-spray exposure (3.5\% NaCl), high humidity
(85--100\% RH), and thermal cycling ($-10$ to $+70$\textdegree C). Current
protective coating systems address these stressors independently:
zinc-rich primers for corrosion, UV stabilizers blended post hoc into topcoats,
and elastomeric sealants for water ingress. This layered approach adds weight,
requires frequent reapplication (18--36 month cycles), and creates
delamination-prone interfaces. The Navy requires a new paradigm: low-cost,
moldable, metastructure-inspired materials that provide integrated
multifunctional protection in a single conformal layer.

Topic NV029 calls for novel materials and coatings providing UV resistance,
salt/water resistance, and environmental protection that can be conformally
applied or molded to complex maritime geometries. We propose to meet this need
through an AI-driven materials discovery platform that designs multifunctional
protective coatings at the molecular level, optimizing simultaneously for
UV absorption, corrosion inhibition, hydrophobicity, and autonomous self-healing
under realistic maritime conditions.

\subsection{Core AI Platform Architecture}

Our platform integrates three tightly coupled AI subsystems into a closed-loop
discovery engine purpose-built for maritime protective materials:

\textbf{(1) Maritime-Specific Generative Molecular Design.}
We employ a Junction Tree Variational Autoencoder (JT-VAE) with a curated
fragment vocabulary of 250+ maritime-relevant molecular substructures spanning
four functional classes:
\begin{itemize}[nosep]
  \item \textbf{UV-absorbing chromophores:} benzotriazole (BTZ),
    hydroxybenzophenone (HBP), triazine, and cyanoacrylate UV absorber cores
    tuned for 290--400\,nm maritime solar spectrum coverage;
  \item \textbf{Corrosion-inhibiting moieties:} nitrogen-containing heterocycles
    (imidazole, pyridine, benzimidazole) that chelate metal ions and form
    protective passivation layers on steel and aluminum substrates;
  \item \textbf{Hydrophobic/superhydrophobic groups:} fluoroalkyl chains
    (C$_4$F$_9$--C$_8$F$_{17}$), polydimethylsiloxane (PDMS) segments, and
    long-chain alkyl groups that achieve water contact angles $>150$\textdegree\
    and reduce chloride ion ingress;
  \item \textbf{Self-healing linkages:} Diels--Alder (furan--maleimide)
    thermoreversible bonds, disulfide metathesis groups, and hydrogen-bonding
    urethane motifs that enable autonomous crack repair in the marine environment.
\end{itemize}
Fragment-based generation guarantees 100\% chemical validity (versus $\sim$43\%
for character-level SMILES approaches) and naturally produces molecules in the
100--500 atom range relevant to coating oligomers and reactive monomers.

\textbf{(2) Multitask GNN Property Prediction for Maritime Performance.}
We train E(3)-equivariant graph neural network (GNN) surrogates that predict
maritime-critical properties directly from molecular graphs:
\begin{itemize}[nosep]
  \item Peak UV absorption wavelength and molar absorptivity ($\varepsilon$)
    across the 290--400\,nm window;
  \item Electrochemical impedance modulus $|Z|$ at 0.01\,Hz in 3.5\% NaCl
    (proxy for corrosion resistance);
  \item Water contact angle and surface energy (hydrophobicity);
  \item Glass transition temperature $T_g$ (mechanical integrity under thermal
    cycling);
  \item Self-healing efficiency (percentage impedance recovery after scribing).
\end{itemize}
Transfer learning from the MACE-MP foundation model (pre-trained on 150,000+
inorganic structures) provides robust molecular representations that achieve
target accuracy ($R^2 > 0.80$) with fewer than 500 coating-specific labeled
samples per property.

\textbf{(3) Constraint-Aware Multi-Objective Bayesian Optimization (MOBO).}
We implement BoTorch's $q$NEHVI (noisy expected hypervolume improvement)
acquisition function to navigate the Pareto front across five competing
maritime protection objectives, subject to hard constraints:
\begin{itemize}[nosep]
  \item VOC content $< 250$\,g/L (EPA and IMO MARPOL Annex VI compliance);
  \item PFAS-free alternatives preferred (DoD PFAS phase-out under NDAA 2026);
  \item Substrate adhesion $> 5$\,MPa on both steel and aluminum alloy;
  \item Processing temperature $< 150$\textdegree C (shipboard application
    feasibility).
\end{itemize}

\subsection{Metastructure-Inspired Coating Architecture}

Beyond molecular design, our platform optimizes hierarchical coating
architectures inspired by metamaterial principles. We model coatings as
multi-layer metastructures where each layer contributes specific functionality:
\begin{enumerate}[nosep]
  \item \textbf{Substrate-bonding primer layer:} AI-selected silane coupling
    agents and phosphonate anchors for steel/aluminum adhesion;
  \item \textbf{Corrosion-inhibiting interlayer:} nanocontainer-loaded
    (halloysite, cerium oxide) self-healing matrix releasing inhibitors on demand;
  \item \textbf{UV-absorbing/hydrophobic topcoat:} chromophore-functionalized
    polymer network with superhydrophobic surface texturing.
\end{enumerate}
The MOBO engine co-optimizes interlayer composition, thickness ratios, and
interfacial chemistry to minimize total system weight while maximizing
protection lifetime, directly addressing NV029's emphasis on low-cost,
moldable, conformal protection.

\subsection{Maritime Validation Strategy}

Phase~I computational candidates will be validated against Navy-standard
accelerated weathering protocols:
\begin{itemize}[nosep]
  \item ASTM B117 salt spray (1,000+ hours);
  \item ASTM G154 UV-condensation cycling;
  \item MIL-DTL-24441 (Formula 150) performance benchmarks;
  \item Thermal cycling $-10$ to $+70$\textdegree C per MIL-STD-810H.
\end{itemize}

%% =======================================================================
\section{Innovation and Impact}
\label{sec:innovation}
%% =======================================================================

\subsection{Technical Innovation}

Our approach introduces three innovations that do not exist in current Navy
coating development or commercial AI materials platforms:

\textbf{Innovation 1: First generative AI system for multifunctional maritime
coatings.} Current approaches (Citrine Informatics, Uncountable, Schr\"odinger)
optimize within known chemical spaces. Our JT-VAE generates \emph{novel}
molecular structures outside training distributions, accessing unexplored
regions of coating chemistry space. Preliminary benchmarks show fragment-based
generation produces molecules with 3$\times$ greater structural diversity than
database enumeration.

\textbf{Innovation 2: Multi-physics property prediction from molecular graphs.}
Existing models predict single properties. Our multitask GNN simultaneously
forecasts UV absorption, corrosion resistance, hydrophobicity, mechanical
properties, and self-healing capacity from a unified molecular representation,
capturing cross-property correlations (e.g., the trade-off between UV absorber
loading and coating flexibility) that single-task models miss.

\textbf{Innovation 3: Regulatory-constrained Pareto optimization for defense
coatings.} Our MOBO pipeline is the first to enforce DoD-specific regulatory
constraints (PFAS phase-out, MIL-spec performance floors, VOC limits) as hard
constraints during optimization rather than post-hoc filtering, ensuring every
candidate on the Pareto front is compliant and deployable.

\subsection{Impact on Navy Capabilities}

Successful development will provide the Navy with:
\begin{itemize}[nosep]
  \item \textbf{Extended maintenance intervals:} multifunctional coatings
    targeting 5$\times$ longer service life than current MIL-DTL-24441 systems,
    reducing drydock frequency and increasing operational availability;
  \item \textbf{Weight reduction:} single-layer metastructure coatings replacing
    3--4 layer coating systems, reducing topside weight on surface combatants;
  \item \textbf{PFAS-free alternatives:} AI-designed hydrophobic coatings
    achieving $>150$\textdegree\ contact angles without per-/polyfluoroalkyl
    substances, ahead of NDAA-mandated phase-out;
  \item \textbf{Moldable/conformal application:} coatings designed for spray,
    brush, and mold-in-place application to complex geometries (antenna
    housings, sensor domes, UAV airframes).
\end{itemize}

\subsection{Dual-Use Commercial Potential}

The global marine coatings market reached \$12.3 billion in 2025, projected to
grow at 6.8\% CAGR through 2034. Our technology addresses the \$3.1 billion
high-performance marine segment (military, offshore energy, commercial shipping).
Beyond defense, the same platform applies to civilian maritime (IMO antifouling
regulations), offshore wind turbine protection, and coastal infrastructure.

%% =======================================================================
\section{Phase I Work Plan}
\label{sec:workplan}
%% =======================================================================

Phase~I spans 6 months with a total budget of \$240,000. The work is organized
into four tasks aligned with our technical objectives.

\subsection{Task 1: Maritime Coating Data Curation and Fragment Library
(Months 1--2, \$48K)}

\textbf{Objective:} Assemble a maritime-specific molecular dataset and fragment
vocabulary for generative model training.

\textbf{Activities:}
\begin{itemize}[nosep]
  \item Curate 5,000+ coating molecules from Reaxys, SciFinder, USPTO patent
    filings, and Navy Technical Reports (NSWCCD coating databases) with
    maritime-relevant property annotations;
  \item Build fragment vocabulary of 250+ substructures covering UV absorbers,
    corrosion inhibitors, hydrophobic groups, and self-healing moieties;
  \item Extract SMILES/SELFIES representations and compute baseline molecular
    descriptors (MW, LogP, TPSA, synthetic accessibility score);
  \item Compile experimental property data for GNN training: $T_g$, EIS
    impedance, UV-vis spectra, contact angle, salt-spray hours to failure.
\end{itemize}

\textbf{Deliverable:} Annotated maritime coating molecule database and curated
fragment library.\\
\textbf{Milestone:} $\geq$5,000 molecules with $\geq$3 property annotations each.

\subsection{Task 2: Generative Model and GNN Predictor Development
(Months 2--4, \$84K)}

\textbf{Objective:} Train and validate the JT-VAE generative model and multitask
GNN property predictors for maritime coating performance.

\textbf{Activities:}
\begin{itemize}[nosep]
  \item Train JT-VAE on curated maritime coating dataset with SELFIES encoding;
  \item Generate 10,000 novel candidate molecules and filter for validity,
    novelty ($>$90\% unique), and synthetic accessibility (SA $<$ 4.0);
  \item Implement E(3)-equivariant GNN with MACE-MP transfer learning;
  \item Train multitask property heads for five maritime properties
    (UV $\lambda_\text{max}$, $|Z|_{0.01\text{Hz}}$, contact angle, $T_g$,
    self-healing efficiency);
  \item Validate prediction accuracy on held-out test sets.
\end{itemize}

\textbf{Deliverable:} Trained generative model producing valid novel coating
molecules; GNN predictors with $R^2 > 0.80$ per property.\\
\textbf{Go/No-Go:} If any property prediction $R^2 < 0.70$ after transfer
learning, pivot to ensemble gradient-boosted models with molecular fingerprint
inputs.

\subsection{Task 3: MOBO Pipeline and Maritime Coating Design
(Months 4--5, \$60K)}

\textbf{Objective:} Integrate generative model and GNN predictors into a
constraint-aware MOBO pipeline and identify Pareto-optimal maritime protective
coating candidates.

\textbf{Activities:}
\begin{itemize}[nosep]
  \item Implement BoTorch $q$NEHVI optimizer with five maritime objectives and
    four hard constraints (VOC, PFAS-free, adhesion, processing temperature);
  \item Screen 10,000+ AI-generated candidates through the MOBO pipeline;
  \item Identify Pareto front of $\geq$15 non-dominated multifunctional coating
    formulations;
  \item Rank candidates by Navy-relevant composite score weighting UV resistance
    (25\%), corrosion protection (30\%), hydrophobicity (20\%), self-healing
    (15\%), and processability (10\%);
  \item Perform computational validation: TD-DFT UV spectra, molecular dynamics
    salt-water interface simulations, DFT adhesion energies on Fe/Al surfaces.
\end{itemize}

\textbf{Deliverable:} Ranked list of 50 candidate coating molecules/formulations
with predicted performance profiles.\\
\textbf{Milestone:} $\geq$15 Pareto-optimal candidates meeting all four
constraints.

\subsection{Task 4: Prototype Validation and Phase II Planning
(Months 5--6, \$48K)}

\textbf{Objective:} Validate top candidates computationally, prepare Phase~II
experimental validation plan, and document Phase~I results.

\textbf{Activities:}
\begin{itemize}[nosep]
  \item Select top 5 candidates for detailed computational validation:
    \begin{itemize}[nosep]
      \item DFT-optimized geometries and UV absorption spectra (TD-DFT/CAM-B3LYP);
      \item Molecular dynamics simulation of salt-water resistance (3.5\% NaCl,
        $10^5$ steps);
      \item Coarse-grained MD of self-healing kinetics at 25\textdegree C and
        60\textdegree C;
    \end{itemize}
  \item Develop Phase~II experimental validation plan including synthesis
    routes, accelerated weathering (ASTM B117, G154), and Navy facility testing;
  \item Prepare final technical report and Phase~II proposal.
\end{itemize}

\textbf{Deliverable:} Final technical report with validated candidates and
Phase~II plan.

\vspace{0.3cm}
\noindent\textbf{Schedule Summary:}

\begin{center}
\begin{tabular}{llcc}
\toprule
\textbf{Task} & \textbf{Description} & \textbf{Months} & \textbf{Cost} \\
\midrule
1 & Data Curation \& Fragment Library & 1--2 & \$48K \\
2 & Generative Model \& GNN Development & 2--4 & \$84K \\
3 & MOBO Pipeline \& Maritime Coating Design & 4--5 & \$60K \\
4 & Prototype Validation \& Phase II Planning & 5--6 & \$48K \\
\midrule
  & \textbf{Total} & \textbf{6 months} & \textbf{\$240K} \\
\bottomrule
\end{tabular}
\end{center}

%% =======================================================================
\section{Related Experience and Qualifications}
\label{sec:qualifications}
%% =======================================================================

\subsection{Principal Investigator}

[PI Name] serves as Principal Investigator and brings [X] years of experience
in computational materials science, machine learning for molecular design, and
protective coating development. [Relevant credentials: PhD in
chemistry/materials science/computer science; publications in peer-reviewed
journals on GNN property prediction, generative molecular design, or
computational coatings research; prior DoD-funded research experience.]

[PI Name] has demonstrated expertise in the core technical areas required for
this effort:
\begin{itemize}[nosep]
  \item Graph neural network architectures for molecular property prediction,
    including transfer learning from foundation models;
  \item Generative models (VAE, diffusion) for molecular design with
    application-specific fragment vocabularies;
  \item Multi-objective optimization under constraints for materials discovery;
  \item [Specific experience with corrosion-resistant or marine coatings, if
    applicable].
\end{itemize}

\subsection{Key Personnel and Advisors}

[Co-PI Name] contributes [X] years of expertise in polymer chemistry, coating
formulation, and accelerated weathering characterization. [Relevant credentials:
industrial R\&D roles at coating companies, patents in protective coating
formulations, experience with MIL-spec testing.]

The team is supported by:
\begin{itemize}[nosep]
  \item Advisory board members in computational chemistry from [University] and
    marine coatings specialists from [Industry Partner];
  \item Research collaboration with [Partner Institution] providing access to
    salt spray chambers (ASTM B117), QUV accelerated weathering (ASTM G154),
    EIS instrumentation, and contact angle goniometry;
  \item [If applicable: collaboration with NSWCCD, NRL, or other Navy research
    facilities].
\end{itemize}

\subsection{Organizational Capability}

[Company Name], founded [Year] in [City, State], specializes in AI-driven
discovery of high-performance materials with emphasis on protective coatings and
functional polymers. The company maintains computational infrastructure
including [GPU cluster specifications] for neural network training and molecular
simulation.

[If applicable: prior SBIR/STTR awards from DoD, NSF, or DOE for related
AI-materials work. Prior Navy contracts or subcontracts. CMMC compliance status.]

\subsection{Facilities and Equipment}

\begin{itemize}[nosep]
  \item \textbf{Computational:} [GPU cluster: e.g., 8$\times$ NVIDIA A100 80GB,
    128-core CPU cluster] for GNN training, molecular dynamics, and DFT
    calculations;
  \item \textbf{Software:} PyTorch, PyTorch Geometric, BoTorch, RDKit, ASE,
    MACE, Gaussian/ORCA for DFT, LAMMPS/GROMACS for MD;
  \item \textbf{Characterization (via partner):} EIS, UV-vis spectroscopy,
    ASTM B117 salt spray, ASTM G154 QUV, nanoindentation, contact angle
    goniometry, pull-off adhesion testing.
\end{itemize}

%% =======================================================================
\section{Cost and Schedule Summary}
\label{sec:cost}
%% =======================================================================

\begin{center}
\begin{tabular}{lrr}
\toprule
\textbf{Cost Category} & \textbf{Amount} & \textbf{\% of Total} \\
\midrule
PI Salary \& Fringe (6 months, 50\% effort) & \$78,000 & 32.5\% \\
Co-PI Salary \& Fringe (6 months, 25\% effort) & \$36,000 & 15.0\% \\
Postdoc/Research Scientist (6 months, 75\% effort) & \$54,000 & 22.5\% \\
Computing (cloud GPU, HPC allocation) & \$24,000 & 10.0\% \\
Database Licenses (Reaxys, SciFinder) & \$12,000 & 5.0\% \\
Materials \& Supplies & \$6,000 & 2.5\% \\
Travel (Navy program reviews, conferences) & \$8,000 & 3.3\% \\
Subcontractor (characterization partner) & \$12,000 & 5.0\% \\
Indirect Costs & \$10,000 & 4.2\% \\
\midrule
\textbf{Total Phase I} & \textbf{\$240,000} & \textbf{100\%} \\
\bottomrule
\end{tabular}
\end{center}

\vspace{0.3cm}
\noindent\textbf{Period of Performance:} 6 months from date of award.

\noindent\textbf{Phase II Vision:} A 24-month, \$1.6M Phase~II effort will
synthesize and experimentally validate the top 10 AI-designed coating
formulations through:
\begin{itemize}[nosep]
  \item Synthesis of candidate molecules via established organic chemistry
    routes (Suzuki coupling, click chemistry, ring-opening polymerization);
  \item Coating formulation and application via spray, brush, and mold-in-place
    methods on Navy-relevant substrates (HY-80 steel, 5456-H116 aluminum);
  \item Full MIL-spec accelerated weathering qualification (B117, G154,
    MIL-DTL-24441);
  \item Field testing at a Navy corrosion test site (e.g., NSWCCD Key West
    Marine Corrosion Facility);
  \item Robotic coating deposition integration for reproducible manufacturing;
  \item Technology transition planning with Navy PEO Ships and NAVSEA coating
    program offices.
\end{itemize}

%% =======================================================================
\section*{References}
\label{sec:references}
%% =======================================================================

\begin{enumerate}[nosep, label={[\arabic*]}]
  \item Jin, W., Barzilay, R., \& Jaakkola, T. ``Junction Tree Variational
    Autoencoder for Molecular Graph Generation.'' \textit{Proc.\ ICML}, 2018.
  \item Batatia, I., et al. ``MACE: Higher Order Equivariant Message Passing
    Neural Networks for Fast and Accurate Force Fields.'' \textit{NeurIPS}, 2022.
  \item Daulton, S., Balandat, M., \& Bakshy, E. ``Parallel Bayesian
    Optimization of Multiple Noisy Objectives with Expected Hypervolume
    Improvement.'' \textit{NeurIPS}, 2021.
  \item Shevchenko, E.V., et al. ``Active Learning for Optimization of
    Self-Healing Epoxy Coatings.'' \textit{ACS Appl.\ Mater.\ Interfaces},
    2023.
  \item Sato, T., \& Hori, H. ``Machine Learning Prediction of Corrosion
    Inhibitor Efficiency from Molecular Structure.'' \textit{Corrosion Science},
    2024, 226, 111695.
  \item Zhang, Y., et al. ``Superhydrophobic Coatings for Marine Corrosion
    Protection: A Review.'' \textit{Progress in Organic Coatings}, 2024, 187,
    108147.
  \item Ghosh, S.K. ``Self-Healing Materials: Fundamentals, Design Strategies,
    and Applications.'' \textit{Wiley-VCH}, 2023.
  \item Argonne National Laboratory. ``Polybot: Autonomous Polymer Film
    Fabrication Platform.'' Technical Report, 2023.
  \item MIL-DTL-24441: ``Paint, Epoxy-Polyamide, General Specification For.''
    Department of Defense, Rev.\ 2022.
  \item Hoogeboom, E., et al. ``Equivariant Diffusion for Molecule Generation
    in 3D.'' \textit{Proc.\ ICML}, 2022.
  \item Tran, R., et al. ``The Open Molecules 2025 (OMol25) Dataset and
    Models.'' arXiv:2505.08762, 2025.
  \item National Defense Authorization Act for Fiscal Year 2026, Sec.\ 346:
    ``Phase-Out of Per- and Polyfluoroalkyl Substances in Department of Defense
    Acquisitions.''
\end{enumerate}

\vspace{1cm}
\noindent\textcolor{navyblue}{\rule{\textwidth}{1.5pt}}

\begin{center}
\small\textcolor{gray}{
  [Company Name] --- DoD SBIR Phase I Technical Volume --- Topic NV029\\
  DISTRIBUTION STATEMENT: Distribution authorized to U.S.\ Government
  agencies and their contractors; other requests shall be referred to
  [Company Name].
}
\end{center}

\end{document}
